\section{Passover, New Exodus, and Covenant}

The Eucharist is intelligible within Israel's worship before it is intelligible
as a theory about time. In Exodus the Passover is simultaneously a dated
deliverance, a household meal, a sacrificial sign, and a perpetual memorial.
Later celebrants do not add another escape from Pharaoh to the first. They
receive their identity from God's accomplished deliverance and submit
themselves anew to its covenantal claim (Exod 12:1--28; 13:3--10). Deuteronomy
can address a later generation as the people with whom the Lord made covenant
at Horeb (Deut 5:2--3). This is corporate and liturgical identification, not a
denial that generations are historically distinct.

The prophets deepen the pattern. Isaiah joins a new exodus, the Servant's
self-offering, covenant renewal, and the feast of salvation (Isa 40:1--11;
42:6; 49:6--12; 52:13--53:12; 55:1--3). Jeremiah promises a new covenant in
which the law is written within and sins are forgiven (Jer 31:31--34).
Ezekiel joins cleansing, a new heart, the Spirit, restored communion, and a
Davidic shepherd (Ezek 36:24--28; 37:21--28). These texts do not constitute a
single hidden script for the upper room. They establish the scriptural world
in which Jesus' covenant cup, commanded memorial, self-gift, and promised
kingdom belong together.

The New Testament then identifies Christ with the paschal fulfillment from
several angles. John the Baptist calls him the Lamb of God (John 1:29, 36).
John's Passion marks paschal preparation and applies the unbroken-bone
scripture to Jesus (John 19:14, 36). Paul calls Christ ``our Pasch'' and joins
the cup and bread to participation in Christ's Blood and Body (1 Cor 5:7--8;
10:16--18). First Peter applies lamb and sacrificial language to redemption
by Christ's precious blood (1 Pet 1:18--21). Revelation presents the slain and
living Lamb as the center of heavenly worship (Rev 5:6--14).

This scriptural complex is simultaneously retrospective, present, and
eschatological. Christian liturgy receives the
fruit of an accomplished exodus, participates in the new covenant now, and
moves toward the Lamb's wedding feast. The Eucharist is memorial because God
has acted; sacrifice because the living High Priest sacramentally makes present
and offers his once-for-all self-offering; communion
because the Church receives and becomes one Body; and pledge because the
kingdom is not yet seen in fullness.

